\documentclass[prl,aps,10pt,superscriptaddress, floatfix]{revtex4-2}
\usepackage[LGR,T1]{fontenc}
\usepackage{amsmath, amssymb, amsthm, mathtools, bbm, physics}
\usepackage{dsfont}
\usepackage{subcaption}
% REVTeX defaults to centered caption, while subcaption assumes justified is the default, so we need to redefine the justified command.
\usepackage{ragged2e} % for the \justifying macro
\DeclareCaptionJustification{justified}{\justifying}

\usepackage[dvipsnames]{xcolor}
\usepackage[page,title, titletoc]{appendix}
\usepackage[colorlinks=true, pdfborder={0 0 0}]{hyperref}
\definecolor{googleblue}{RGB}{34, 0, 204}
\definecolor{panblue}{RGB}{0,24,150}
\definecolor{carmine}{RGB}{150, 0, 24}
\hypersetup{
unicode=true,
bookmarksnumbered=false,
bookmarksopen=false,
breaklinks=false,
colorlinks=true,
linkcolor=carmine,
citecolor=googleblue,
urlcolor=panblue,
anchorcolor=OliveGreen}

\newcommand{\tmp}[1]{\textcolor{red}{#1}}
\newcommand{\pedro}[1]{{\color{purple} #1}}
\newcommand{\bereket}[1]{{\color{blue} #1}}
\newcommand{\elie}[1]{{\color{RawSienna} #1}}

\usepackage{newtxtext}

%\usepackage{microtype}
%\microtypecontext{spacing=nonfrench}
%\microtypesetup{
%expansion=false,
%protrusion={true,nocompatibility},
%activate={true,nocompatibility},
%tracking=false,
%kerning=true,
%spacing={true}
%}

\usepackage[all=normal,floats=tight,mathspacing=tight,wordspacing=tight,paragraphs=normal,tracking=tight,charwidths=tight,mathdisplays=normal,sections=normal,margins=normal]{savetrees}

%\usepackage[backend=biber]{biblatex}
%\addbibresource{main.bib}

\usepackage{tikz} % TikZ
\usetikzlibrary{calc, arrows, patterns, shapes, decorations, arrows.meta, fit, positioning}
\tikzset{
    auto, node distance =1 cm and 1 cm,semithick,
    var/.style ={circle, draw, minimum width = 1cm, ultra thick},
    latent/.style ={regular polygon, regular polygon sides=3, inner sep=1pt, draw, minimum width = 1.2cm, ultra thick},
    point/.style = {circle, draw, inner sep=0.06cm, fill, node contents={}},
    triangle/.style = {regular polygon, regular polygon sides=3, draw, inner sep=0.06cm, fill, node contents={}},
    bidir/.style={Latex-Latex,dashed},
    dir/.style={-Latex, thick},
    el/.style = {inner sep=2pt, align=left, sloped}
}
\tikzstyle{vertex}=[circle, fill=black!10, draw=black]
\tikzstyle{edge}=[thick]
\tikzstyle{clique}=[line width=4, draw=black!70]

% Theorem environments
\newtheorem{theo}{Theorem}
\newtheorem{lemma}{Lemma}
\newtheorem{prop}{Proposition}
\newtheorem{coro}{Corollary}
\newtheorem*{defi}{Definition}
\newtheorem{obs}{Observation}

% Useful math shortcuts
\newcommand{\eye}{\mathbbm{1}}
\newcommand{\mO}{\mathcal{O}}
\newcommand{\mL}{\mathcal{L}}
\newcommand{\mH}{\mathcal{H}}
\newcommand{\Cset}{\mathcal{C}}
\newcommand{\Qset}{\mathcal{Q}}
\newcommand{\Nset}{\mathcal{N}}
\newcommand{\NBset}{\mathcal{N}_\mathrm{B}}
\newcommand{\Real}{\mathbb{R}}
\newcommand{\Nat}{\mathbb{N}}
\graphicspath{{imgs/}}

\DeclareMathOperator{\Pa}{Pa}

\DeclareMathOperator{\pa}{pa}
\DeclareMathOperator{\Ch}{Ch}
\DeclareMathOperator{\opa}{opa}
\DeclareMathOperator{\OPa}{OPa}
\DeclareMathOperator{\LPa}{LPa}

\usepackage{tikz}
% drawing parameters
\newcommand{\circdist}{1}  % distance from origin to center of circles
\newcommand{\circrad}{6/4} % radius of the circles
\newcommand{\circlethickness}{2.5mm} % uh, thickness of the circles

% distance from the origin to the three "interior" intersections
\pgfmathsetmacro{\intrad}{sqrt((\circrad)^2 - 3*(\circdist)^2/4) - \circdist/2}

% distance from the origin to the three "exterior" intersections
\pgfmathsetmacro{\extrad}{sqrt((\circrad)^2 - 3*(\circdist)^2/4) + \circdist/2}

% so we can just specify an angle and get the correct color for the circle
\colorlet{180}{blue!60}
\colorlet{60}{red!50}
\colorlet{300}{green!40}

% draws one of our circles
\newcommand{\mycircle}[1]{%
  \draw[thick, double distance=\circlethickness, double=#1]
  (#1:\circdist) circle (\circrad);}


%%% Adrian Commands %%%%%
\usepackage{xcolor}
\definecolor{pansypurple}{rgb}{0.47, 0.09, 0.29}

\newcommand\al[1]{\textcolor{pansypurple}{ [A:\,#1]}}
\newcommand\p[1]{\textcolor{blue}{ [P:\,#1]}}
%%%%%%%%%%%%%%%%%%%%%%%%%%

\begin{document}
\title{Generalized Probabilistic Topological Field Theories: Notes}
\author{Pedro Lauand}
\email{plauand@perimeterinstitute.ca}
\affiliation{Perimeter Institute for Theoretical Physics, Waterloo, Ontario, Canada, N2L 2Y5}

\maketitle
%\onecolumn
\section{Introduction}

\al{Eventually, I think the main theme of the introduction of an eventual paper should be around the notion of locality. At some point high-energy theorists realized that QFTs are things you ``mount" on mafilods, or ``live" on them. Locality comes when you study the compatibility of gluing: if you have two processes labelled by two manifolds, the process associated to the gluing (when admissible in the sense that output mfld of one equals input one of the other), is assigned the process of the composition. This compatibility of composition, means that there is no non-local nature to time evolution, in the sense that combining processes would require some new notion of composition that would need to somehow know more about everything that happened in the joint process. Sorry for the convoluted sentence, I can explain it better in person. In any case, there should be similar statements about spatial locality and extended TFTs.}
\\
\al{Another important theme might be what exists beyond QFTs}

In the functorial approach, an $n$-dimensional topological quantum field theory (TQFT) is specified by a symmetric monoidal functor
\[
Z: \mathbf{Bord}_n \rightarrow \mathcal{C},
\]
where $\mathbf{Bord}_n$ is the $n$-dimensional bordism category and $\mathcal{C}$ is typically a symmetric monoidal category such as $\mathbf{Vec}$ or $\mathbf{Hilb}$ \al{Idk if people actually make it valued in \textbf{Hilb}}. The category $\mathbf{Bord}_n$ has closed $(n-1)$-manifolds as objects and $n$-dimensional bordisms as morphisms; composition is given by gluing, while disjoint union defines the monoidal product. In this formulation, the choice of $\mathcal{C}$ is part of the structure of the theory itself: it determines the nature of the objects assigned to hypersurfaces and of the processes assigned to bordisms. From this perspective, one may regard a TQFT as providing a rule by which topological spacetime data are organized \al{Idk how I feel about this perspective just yet}.

The traditional choice of target category in functorial TQFT is motivated by quantum theory itself.  Indeed, taking $\mathcal {C}$ to be \textbf{Vec}, or \textbf{Hilb} , reflects the idea that a hypersurface \al{time slice might be clearer} should determine a space of states and that a bordism should determine an evolution between such spaces. In this way, the functorial formalism extracts from quantum theory the linear structures attached to states and processes, and packages them in a purely topological setting. From the empirical point of view, however, quantum theory is a theory of probabilities: its physical content is accessed through the probabilities it assigns to possible measurement outcomes. Accordingly, while TQFTs capture the linear structures underlying quantum theory, it does not by itself provide a direct probabilistic interpretation of the associated spacetime data.  To achieve this, one must instead consider a choice of target category whose structure encodes operational probabilistic meaning \al{It's more subtle than this. TFTs are usually Euclidean QFTs, in which case the notion of uintarity (and therefore of probablity) is encoded in a property called \textit{reflexion positivity}.}.

In recent decades, the study of the foundations of quantum theory has led to the development of broader operational frameworks aimed at isolating its essential physical features. Among these, generalized probabilistic theories (GPTs) provide a setting that captures the minimal structure required to describe physical theories in probabilistic terms. Within this framework, both classical probability theory and quantum theory arise as particular cases, while more general  probabilistic models are also permitted. As such, GPTs offer a unified language in which the probabilistic content of physical theories can be formulated.

Any GPT admits a natural categorical formulation, in which physical systems are regarded as objects and allowed physical transformations as morphisms, with composition corresponding to sequential processes and a monoidal structure encoding parallel composition \al{Something that cliked for me in our conversations, is that the monoidal structure is just how you describe a joint system; you need a notion of what happens when you have system A \textit{and} system B}. In this way, GPTs give rise to symmetric monoidal categories of physical processes. Motivated by this observation, we will refer to topological quantum field theories valued in such categories as generalized probabilistic topological field theories, or GPTFTs.

\section{Category of GPTs}

A generalized probabilistic theory seeks to capture the minimal ingredients of a physical theory from a purely operational point of view. In this setting, physical systems and their allowed transformations are organized into a symmetric monoidal category, subject to additional requirements reflecting the probabilistic nature of the theory. The two basic ingredients of any such theory are \emph{states}, which encode preparation procedures, and \emph{effects}, which encode measurement outcomes; these are treated as co-equal primitives, coupled by a probability rule. We work throughout with finite-dimensional real vector spaces.

\subsection{Objects: systems}

To each physical system $A$ we associate two finite-dimensional real vector spaces: an ambient \emph{state space} $V_A$ and an ambient \emph{effect space} $W_A$. The operationally meaningful states and effects are not all of $V_A$ and $W_A$, but only those lying in distinguished cones, which we now introduce.

Recall that a subset $C$ of a real vector space $V$ is a \emph{proper cone} if it is
\begin{itemize}
\item[(a)] \emph{closed} (as a subset of $V$ in the Euclidean topology);
\item[(b)] \emph{convex} and stable under nonnegative scaling, i.e. $\lambda x + \mu y \in C$ for all $x,y\in C$ and $\lambda,\mu\ge 0$;
\item[(c)] \emph{pointed}, $C\cap(-C)=\{0\}$ (it contains no line through the origin);
\item[(d)] \emph{generating}, $C-C=V$ (its linear span is all of $V$).
\end{itemize}
Properness is exactly what is needed for a cone to play the role of a set of physical configurations: closedness ensures that limits of physical states are again physical, convexity encodes the possibility of probabilistic mixing, pointedness rules out a configuration and its negative both being physical, and the generating property ensures that the ambient space contains no directions unrelated to physics.

The data of a system $A$ are then as follows.
\begin{itemize}
\item A \emph{state cone} $V_A^+\subseteq V_A$, a proper cone whose elements are the (unnormalized) states. Since $V_A^+$ is generating, $V_A=\operatorname{span}(V_A^+)$, so the ambient state space is recovered from the cone by linearization.
\item An \emph{effect cone} $W_A^+\subseteq W_A$, a proper cone whose elements are the effects, together with a distinguished \emph{unit effect} $u_A\in W_A^+$. Again $W_A=\operatorname{span}(W_A^+)$.
\item A \emph{probability rule}: a bilinear pairing
\[
B_A : V_A \times W_A \to \mathbb{R},
\]
whose value $B_A(v,w)$ is interpreted, for a (normalized) state $v$ and an effect $w$, as the probability that the measurement outcome $w$ occurs given the preparation $v$.
\end{itemize}

The pairing is required to satisfy three conditions.
\begin{itemize}
\item[(i)] \emph{Positivity on the cones:}
\[
B_A(V_A^+,W_A^+)\subseteq \mathbb{R}^{\ge 0},
\]
i.e. $B_A(v,w)\ge 0$ for all $v\in V_A^+$ and $w\in W_A^+$. Operationally, probabilities are nonnegative, and this must hold on every state--effect pair that the theory deems physical.
\item[(ii)] \emph{Separation in each argument:} $B_A$ is nondegenerate separately in $v$ and in $w$. That is, if $B_A(v,w)=B_A(v',w)$ for all $w\in W_A$, then $v=v'$; and if $B_A(v,w)=B_A(v,w')$ for all $v\in V_A$, then $w=w'$. Operationally, two states that assign the same probability to every effect are physically indistinguishable and are therefore identified, and dually for effects. This condition guarantees that $V_A$ and $W_A$ carry no redundant degrees of freedom beyond those that are operationally accessible through the probability rule.
\item[(iii)] \emph{Strict-positive normalization:}
\[
B_A(v,u_A)>0 \qquad \text{for all } v\in V_A^+\setminus\{0\}.
\]
Operationally, $u_A$ is the trivial measurement outcome that always occurs, so $B_A(v,u_A)$ is the total weight (the success probability) of the unnormalized preparation $v$. Condition (iii) states that every nonzero physical preparation has strictly positive total weight: a state with zero total weight would never be prepared at all and so does not correspond to a genuine, nonzero physical state. Geometrically, $u_A$ defines a hyperplane $\{v: B_A(v,u_A)=1\}$ that meets the cone $V_A^+$ transversally and bounds a compact slice.
\end{itemize}

\paragraph{Derived structure.}
The remaining operational notions are not posited separately: they are \emph{derived} from the data $(V_A^+,W_A^+,u_A,B_A)$ above.
\begin{itemize}
\item The \emph{normalized states} are the slice of the state cone selected by the unit effect,
\[
\Omega_A:=\{\omega\in V_A^+: B_A(\omega,u_A)=1\}.
\]
By condition (iii) this is a base of the cone $V_A^+$: every nonzero $v\in V_A^+$ is uniquely written as $v=B_A(v,u_A)\,\omega$ with $\omega\in\Omega_A$. We assume $\Omega_A$ is convex and compact; operationally, this means that if a sequence of preparations $(\omega_n)$ is such that $B_A(\omega_n,e)$ converges for every effect $e$, then the limiting probability assignment is realized by some $\omega\in\Omega_A$. Starting from a normalized preparation $\omega\in\Omega_A$, a procedure that succeeds only with probability $p$ is represented by the subnormalized state $p\,\omega\in V_A^+$; thus the cone $V_A^+$ already contains both normalized and subnormalized states.
\item The \emph{allowed effects} are those effects bounded above by the unit,
\[
\mathcal E_A:=[0,u_A]=\{e\in W_A^+: u_A-e\in W_A^+\}.
\]
For a normalized state $\omega\in\Omega_A$ and an allowed effect $e\in\mathcal E_A$, positivity together with $e\le u_A$ gives $B_A(\omega,e)\in[0,1]$, as a probability must.
\item A \emph{measurement} is a finite family of effects $(e_i)_i$ summing to the unit,
\[
\sum_i e_i = u_A.
\]
Each $e_i$ is a possible outcome, and the normalization expresses that exactly one outcome occurs: for a normalized state $\omega$ the outcome probabilities $P_A(i):=B_A(\omega,e_i)$ are nonnegative and sum to $B_A(\omega,u_A)=1$.
\end{itemize}
We note in passing that the bilinearity of $B_A$ is precisely the statement that the probability rule respects convex mixing of preparations and of measurements: $B_A(p\,\omega_1+(1-p)\,\omega_2,e)=p\,B_A(\omega_1,e)+(1-p)\,B_A(\omega_2,e)$, and similarly in the second argument. Operationally, randomizing which of two preparation (or measurement) procedures is carried out is faithfully represented at the level of probabilities.

\subsection{Morphisms: transformations}

A transformation describes a physical process taking one system into another. Operationally, such a process must do two things at once: it must turn preparations of $A$ into preparations of $B$ (acting on states, in the forward direction), and it must turn measurements on $B$ into measurements on $A$ (acting on effects, in the backward direction), in a way that is consistent with the probability rule. Accordingly, a \emph{transformation} $\Gamma:A\to B$ is a pair of linear maps
\[
\Gamma_\Omega : V_A \to V_B, \qquad \Gamma^{\mathcal E} : W_B \to W_A,
\]
related by the \emph{adjunction}
\begin{equation}\label{eq:adjunction}
B_B(\Gamma_\Omega v,\,w) = B_A(v,\,\Gamma^{\mathcal E} w)
\qquad \forall\,(v,w)\in V_A\times W_B.
\end{equation}
Equation~\eqref{eq:adjunction} is the operational consistency requirement: the probability of obtaining outcome $w$ on the transformed state $\Gamma_\Omega v$ must equal the probability of obtaining the pulled-back outcome $\Gamma^{\mathcal E} w$ on the original state $v$. It is immaterial whether one pushes the state forward and measures, or pulls the effect back and measures the original preparation.

The two maps are further required to preserve the respective primitives, which is expressed by the two positivity conditions
\begin{equation}\label{eq:Gamma-positivity}
\Gamma_\Omega(V_A^+)\subseteq V_B^+,
\qquad
\Gamma^{\mathcal E}(W_B^+)\subseteq W_A^+.
\end{equation}
The first sends valid states to valid states; the second sends valid effects to valid effects.

Finally, a transformation may be \emph{probabilistic}: it need not succeed deterministically, so applying it to a normalized state may produce a subnormalized one. This is encoded by requiring that the total weight never increase,
\[
B_B(\Gamma_\Omega\omega,\,u_B)\le B_A(\omega,\,u_A)\qquad \forall\,\omega\in\Omega_A,
\]
which, using the adjunction~\eqref{eq:adjunction} with $w=u_B$ and the fact that $\Omega_A$ spans $V_A^+$, is equivalent to the single condition
\begin{equation}\label{eq:subunital}
\Gamma^{\mathcal E}(u_B)\le u_A.
\end{equation}
We call a transformation satisfying~\eqref{eq:subunital} \emph{operationally allowed} (probability-non-increasing). For the purposes of GPTFTs it is convenient to retain the larger, cone-linear class of morphisms defined only by the adjunction~\eqref{eq:adjunction} and the positivity~\eqref{eq:Gamma-positivity}, with the operationally allowed maps~\eqref{eq:subunital} singled out as a distinguished subclass. In this larger category, $\operatorname{Hom}(A,B)$ is itself a cone rather than a vector space.

Composition is read off from the two directions in which the maps act: for $\Gamma:A\to B$ and $\Delta:B\to C$,
\[
(\Delta\circ\Gamma)_\Omega := \Delta_\Omega\circ\Gamma_\Omega,
\qquad
(\Delta\circ\Gamma)^{\mathcal E} := \Gamma^{\mathcal E}\circ\Delta^{\mathcal E},
\]
the effect side composing in the opposite order, as the adjunction~\eqref{eq:adjunction} demands. The identity transformation on $A$ is the pair $(\operatorname{id}_{V_A},\operatorname{id}_{W_A})$, which trivially satisfies~\eqref{eq:adjunction}--\eqref{eq:subunital}.

\subsection{Monoidal structure: composite systems}

To describe composite systems we introduce a monoidal product $\,\square\,$. Operationally, $A\square B$ is the joint system obtained by regarding $A$ and $B$ as independent subsystems; the monoidal structure is simply the rule for what it means to have system $A$ \emph{and} system $B$ at once. Composition must be compatible with the probabilistic structure, in that one can form product states and product effects whose joint probabilities factorize. Concretely, $\square$ supplies bilinear product maps on states and on effects,
\[
\square : V_A \times V_B \to V_{A\square B},
\qquad
\square : W_A \times W_B \to W_{A\square B},
\]
subject to the \emph{factorizing pairing}
\begin{equation}\label{eq:factorizing}
B_{A\square B}(v\square v',\, w\square w') = B_A(v,w)\,B_B(v',w')
\end{equation}
for all $v\in V_A$, $v'\in V_B$, $w\in W_A$, $w'\in W_B$. Operationally, on product preparations measured by product effects the joint probability is the product of the marginal probabilities.

There is a distinguished system $\mathbf{1}$, the \emph{monoidal unit}, representing the trivial system, characterized by $A\square\mathbf 1\cong A\cong\mathbf 1\square A$. The trivial system has a single normalized state and a single deterministic effect, so at the cone level
\[
V_{\mathbf 1}^+\cong\mathbb{R}^{\ge 0},
\qquad
W_{\mathbf 1}^+\cong\mathbb{R}^{\ge 0},
\qquad\text{whence}\qquad
\mathbf 1=(\mathbb{R},\mathbb{R}^{\ge 0}).
\]

\paragraph{The composite ambient space.}
We now extract the generic structure of the composite spaces $V_{A\square B}$ and $W_{A\square B}$. By the universal property of the tensor product, the bilinear map $\square:V_A\times V_B\to V_{A\square B}$ factors through a unique linear map
\[
\phi_\square : V_A\otimes V_B \to V_{A\square B},
\qquad
\phi_\square(v\otimes v')=v\square v'.
\]
We claim that the operational requirements force $\phi_\square$ to be \emph{injective}.

To prove this, introduce the bilinear map $\widetilde B_{A\otimes B}:(V_A\otimes V_B)\times(W_A\otimes W_B)\to\mathbb{R}$ defined on simple tensors by
\[
\widetilde B_{A\otimes B}(v\otimes v',\,w\otimes w')=B_A(v,w)\,B_B(v',w').
\]
We first record that $\widetilde B_{A\otimes B}$ is separating in each argument, inherited from the separation of $B_A$ and $B_B$. Indeed, let $x\in V_A\otimes V_B$ satisfy $\widetilde B_{A\otimes B}(x,\,w\otimes w')=0$ for all $w\in W_A$ and $w'\in W_B$. Writing $x=\sum_i v_i\otimes v_i'$ with the $v_i\in V_A$ linearly independent, we compute, for every fixed $w'$,
\[
\widetilde B_{A\otimes B}(x,\,w\otimes w')
=\sum_i B_A(v_i,w)\,B_B(v_i',w')
=B_A\!\Big(\sum_i B_B(v_i',w')\,v_i,\;w\Big)=0
\quad\forall w.
\]
Separation of $B_A$ in its first argument gives $\sum_i B_B(v_i',w')\,v_i=0$, and linear independence of the $v_i$ then forces $B_B(v_i',w')=0$ for every $i$. As this holds for all $w'$, separation of $B_B$ yields $v_i'=0$ for every $i$, hence $x=0$. The same argument with the roles exchanged gives separation in the second argument.

Now the factorizing pairing~\eqref{eq:factorizing} says exactly that, on simple tensors, $B_{A\square B}\circ(\phi_\square\times(\,\cdot\,\square\,\cdot\,))$ reproduces $\widetilde B_{A\otimes B}$; by bilinearity this extends to
\begin{equation}\label{eq:compat-pairing}
B_{A\square B}\big(\phi_\square(x),\,w\square w'\big)=\widetilde B_{A\otimes B}(x,\,w\otimes w')
\qquad\forall x\in V_A\otimes V_B,\;\forall w,w'.
\end{equation}
Suppose $\phi_\square(x)=0$. Then the left-hand side of~\eqref{eq:compat-pairing} vanishes for all $w,w'$, so $\widetilde B_{A\otimes B}(x,\,w\otimes w')=0$ for all $w,w'$, and separation of $\widetilde B_{A\otimes B}$ forces $x=0$. Hence $\ker\phi_\square=\{0\}$ and $\phi_\square$ is injective, as claimed.

Injectivity lets us identify $\operatorname{Im}(\phi_\square)$ with $V_A\otimes V_B$, so the composite ambient state space splits as
\begin{equation}\label{eq:V-decomp}
V_{A\square B}\cong (V_A\otimes V_B)\oplus C,
\end{equation}
where $C$ is a complement to $\operatorname{Im}(\phi_\square)$. An entirely analogous argument on the effect side gives
\begin{equation}\label{eq:W-decomp}
W_{A\square B}\cong (W_A\otimes W_B)\oplus D.
\end{equation}
We call $C$ and $D$ the \emph{global sectors}: they account for degrees of freedom of the joint system that are not accessible to either subsystem alone, lying outside the locally generated tensor-product sector. The vanishing of the global sector is exactly the property of local tomography:
\[
C=0
\quad\Longleftrightarrow\quad
\text{the theory is locally tomographic,}
\]
i.e. every joint state is completely determined by the joint statistics of local measurements (and dually $D=0$ on the effect side). In general the global sectors need not vanish; keeping them is essential for theories, such as fermionic ones, that genuinely fail local tomography.

\paragraph{Underdetermination of the composite cone.}
The decomposition~\eqref{eq:V-decomp} fixes the composite ambient \emph{space} but not the composite \emph{cone} $V_{A\square B}^+$. Operationally, two extreme choices are always available. The \emph{minimal} cone $\otimes_{\min}$ is generated by the product states $v\square v'$ with $v\in V_A^+$, $v'\in V_B^+$, i.e. the convex closure of separable preparations. The \emph{maximal} cone $\otimes_{\max}$ consists of all elements assigned nonnegative probability by every product effect, i.e. those $g$ with $B_{A\square B}(g,\,w\square w')\ge 0$ for all $w\in W_A^+$, $w'\in W_B^+$. Compatibility with the local probabilistic structure forces
\begin{equation}\label{eq:cone-underdet}
\otimes_{\min}\ \subseteq\ V_{A\square B}^+\ \subseteq\ \otimes_{\max},
\end{equation}
but leaves the precise composite cone \emph{underdetermined}: any proper cone lying between these bounds (and compatible with the global sector) is an admissible choice of composite. The freedom in~\eqref{eq:cone-underdet}, together with the possibly nonzero global sectors $C,D$, is precisely where a GPT can depart from a theory determined by its subsystems.

\subsection{The category of GPTs}

Collecting the data above, and passing to the linearized (cone-linear) setting, we can package the whole construction as a single definition.

\begin{defi}[Category of GPTs]
An \emph{object} (system $A$) consists of finite-dimensional real vector spaces $V_A$ and $W_A$, proper cones $V_A^+\subseteq V_A$ and $W_A^+\subseteq W_A$ with $V_A=\operatorname{span}(V_A^+)$ and $W_A=\operatorname{span}(W_A^+)$, a distinguished unit effect $u_A\in W_A^+$, and a bilinear pairing $B_A:V_A\times W_A\to\mathbb{R}$ that is (i) positive on the cones, $B_A(V_A^+,W_A^+)\subseteq\mathbb{R}^{\ge0}$, (ii) separating in each argument, and (iii) strictly normalized, $B_A(v,u_A)>0$ for all $v\in V_A^+\setminus\{0\}$. The normalized states and allowed effects are the derived objects
\[
\Omega_A=\{\omega\in V_A^+:B_A(\omega,u_A)=1\},
\qquad
\mathcal E_A=[0,u_A]=\{e\in W_A^+: u_A-e\in W_A^+\}.
\]

A \emph{morphism} $\Gamma:A\to B$ is a pair of linear maps $(\Gamma_\Omega,\Gamma^{\mathcal E})$ with $\Gamma_\Omega:V_A\to V_B$ and $\Gamma^{\mathcal E}:W_B\to W_A$ satisfying the adjunction $B_B(\Gamma_\Omega v,w)=B_A(v,\Gamma^{\mathcal E}w)$ together with the positivity conditions $\Gamma_\Omega(V_A^+)\subseteq V_B^+$ and $\Gamma^{\mathcal E}(W_B^+)\subseteq W_A^+$; composition is $(\Delta\circ\Gamma)_\Omega=\Delta_\Omega\circ\Gamma_\Omega$, $(\Delta\circ\Gamma)^{\mathcal E}=\Gamma^{\mathcal E}\circ\Delta^{\mathcal E}$, and the identity is $(\operatorname{id},\operatorname{id})$. The \emph{operationally allowed} transformations are the subclass additionally satisfying $\Gamma^{\mathcal E}(u_B)\le u_A$.

The category is symmetric monoidal under $\square$, with unit $\mathbf 1=(\mathbb{R},\mathbb{R}^{\ge0})$: the composite $A\square B$ carries bilinear product maps on states and effects whose pairing factorizes, $B_{A\square B}(v\square v',w\square w')=B_A(v,w)\,B_B(v',w')$, the ambient spaces decompose as $V_{A\square B}\cong(V_A\otimes V_B)\oplus C$ and $W_{A\square B}\cong(W_A\otimes W_B)\oplus D$ with global sectors $C,D$ (vanishing iff the theory is locally tomographic), and the composite cone is any admissible choice with $\otimes_{\min}\subseteq V_{A\square B}^+\subseteq\otimes_{\max}$.
\end{defi}

\subsection{From operational data to Vec-like systems}

The data assembled above are deliberately redundant: states and effects are treated as co-equal primitives, a system carries two ambient spaces $(V_A,W_A)$ coupled by an abstract pairing $B_A$, and a transformation is a \emph{pair} of maps $(\Gamma_\Omega,\Gamma^{\mathcal E})$. In finite dimension this redundancy can be removed entirely. The key is the separation requirement on $B_A$: because the pairing has no kernel in either argument, every effect is already determined by the numbers it assigns to states, and every transformation on states already fixes its action on effects. We now make this precise through two collapses, each a direct consequence of separation, after which a system is described purely by a vector space, its state cone, and an effect cone living in the dual space --- data internal to $\mathbf{Vec}$.

\subsubsection*{Effects are functionals and the pairing is evaluation}

Recall the canonical evaluation pairing $\langle\cdot,\cdot\rangle:V_A\times V_A^*\to\mathbb{R}$, $\langle v,\varphi\rangle=\varphi(v)$, between $V_A$ and its dual space $V_A^*$. The abstract effect space $(W_A,B_A)$ maps into this dual via the \emph{flat map}
\[
\flat_A:W_A\longrightarrow V_A^*,\qquad \flat_A(w):=B_A(-,w),
\]
which sends an effect $w$ to the linear functional $v\mapsto B_A(v,w)$. The map $\flat_A$ is linear because $B_A$ is bilinear, and it is built precisely so that it transports $B_A$ to the canonical evaluation:
\begin{equation}\label{eq:flat-intertwines}
\langle v,\flat_A(w)\rangle=\flat_A(w)(v)=B_A(v,w)\qquad\forall\, v\in V_A,\ w\in W_A.
\end{equation}

\begin{lemma}[Effects are functionals]\label{lem:flat}
Let $(V_A,W_A,B_A)$ be a finite-dimensional system whose pairing $B_A$ separates states and effects. Then $\flat_A:W_A\to V_A^*$ is a linear isomorphism; in particular $\dim W_A=\dim V_A$. Under $\flat_A$ the pairing $B_A$ becomes the canonical evaluation $\langle\cdot,\cdot\rangle$, in the sense of \eqref{eq:flat-intertwines}.
\end{lemma}

\begin{proof}
We show $\flat_A$ is injective and surjective.

\emph{Injectivity (from separation in the effect argument).} Suppose $\flat_A(w)=0$. By definition this means $B_A(v,w)=0$ for every $v\in V_A$, i.e.\ the effect $w$ assigns the same value, zero, as the effect $0$ does on all states. Separation in the effect argument identifies effects that agree on all states, hence $w=0$. Thus $\ker\flat_A=\{0\}$ and $\flat_A$ is injective.

\emph{Surjectivity (by a dimension count using separation in the state argument).} Injectivity already gives $\dim W_A\le\dim V_A^*=\dim V_A$. For the reverse inequality, consider the symmetric construction
\[
\sharp_A:V_A\longrightarrow W_A^*,\qquad \sharp_A(v):=B_A(v,-),
\]
which is injective by separation in the \emph{state} argument, exactly as above with the roles exchanged. Hence $\dim V_A\le\dim W_A^*=\dim W_A$. Combining the two inequalities forces $\dim W_A=\dim V_A$. An injective linear map between spaces of equal finite dimension is an isomorphism, so $\flat_A$ is surjective as well.

Finally, \eqref{eq:flat-intertwines} holds by the very definition of $\flat_A$, so the isomorphism carries $B_A$ to the canonical evaluation.
\end{proof}

Lemma~\ref{lem:flat} lets us discard the abstract effect data altogether. We henceforth identify
\[
W_A\;\cong\;V_A^*,\qquad B_A(v,w)\;=\;\langle v,w\rangle\;=\;w(v),
\]
so that effects are simply linear functionals on states and the probability rule is their evaluation. The positivity condition $B_A(V_A^+,W_A^+)\subseteq\mathbb{R}^{\ge0}$ becomes the statement that every positive effect is nonnegative on every positive state; equivalently
\begin{equation}\label{eq:Wpos-in-dual}
W_A^+\;\subseteq\;(V_A^+)^*=\{\varphi\in V_A^*:\varphi(v)\ge0\ \forall v\in V_A^+\},
\end{equation}
the effect cone is contained in the dual cone of the state cone.

It is worth isolating exactly how much freedom remains in $W_A^+$. The \emph{no-restriction hypothesis} is the assertion that \emph{all} mathematically admissible effects are physically allowed, i.e.\ that \eqref{eq:Wpos-in-dual} is an equality,
\[
\text{(no restriction)}\qquad W_A^+=(V_A^+)^*.
\]
We do \emph{not} assume it. Dropping it leaves $W_A^+$ a proper subcone of $(V_A^+)^*$, and the operational axioms then constrain this subcone in exactly two ways:
\begin{itemize}
\item[(a)] $W_A^+$ \emph{spans} $V_A^*$, i.e.\ $\operatorname{span}(W_A^+)=V_A^*$. This is precisely separation in the state argument re-read through $\flat_A$: if the effects spanned only a proper subspace of $V_A^*$, some nonzero $v\in V_A$ would be annihilated by every effect and hence indistinguishable from $0$. Operationally, \emph{effects separate states}. We will see below that this same spanning property is what guarantees the dual object's state cone is generating.
\item[(b)] $u_A\in W_A^+$, the unit effect is a legitimate positive effect.
\end{itemize}
These two conditions --- spanning and containment of the unit --- are the \emph{only} residue of refusing the no-restriction hypothesis. Everything else about the effect side is now subsumed into the geometry of the single cone $W_A^+\subseteq V_A^*$.

\subsubsection*{A transformation is a single map}

With effects realized as functionals, the two halves of a transformation are no longer independent. Recall that the transpose of a linear map $f:V_A\to V_B$ is the map $f^*:V_B^*\to V_A^*$ defined by $\langle v,f^*(\varphi)\rangle=\langle f(v),\varphi\rangle$ for all $v\in V_A$, $\varphi\in V_B^*$.

\begin{lemma}[A transformation is one map]\label{lem:onemap}
Let $\Gamma=(\Gamma_\Omega,\Gamma^{\mathcal E})$ be a transformation $A\to B$, with $\Gamma_\Omega:V_A\to V_B$ and $\Gamma^{\mathcal E}:W_B\to W_A$ satisfying the adjunction $B_B(\Gamma_\Omega v,w)=B_A(v,\Gamma^{\mathcal E}w)$. Under the identifications $W_A\cong V_A^*$, $W_B\cong V_B^*$ of Lemma~\ref{lem:flat}, one has
\[
\Gamma^{\mathcal E}=(\Gamma_\Omega)^*.
\]
The effect-side map is therefore not independent data: it is the transpose of the state-side map.
\end{lemma}

\begin{proof}
Write $f:=\Gamma_\Omega$ and identify effects with functionals via the flat maps, so that pairings become evaluations. The adjunction reads, for all $v\in V_A$ and $w\in W_B\cong V_B^*$,
\[
\langle f(v),w\rangle=B_B(\Gamma_\Omega v,w)=B_A(v,\Gamma^{\mathcal E}w)=\langle v,\Gamma^{\mathcal E}w\rangle.
\]
By the defining property of the transpose, $\langle f(v),w\rangle=\langle v,f^*(w)\rangle$ as well. Hence $\langle v,\Gamma^{\mathcal E}w\rangle=\langle v,f^*(w)\rangle$ for all $v\in V_A$. Separation in the state argument (equivalently, nondegeneracy of $\langle\cdot,\cdot\rangle$) gives $\Gamma^{\mathcal E}w=f^*(w)$ for every $w$, i.e.\ $\Gamma^{\mathcal E}=f^*=(\Gamma_\Omega)^*$.
\end{proof}

Consequently a morphism is a \emph{single} positive linear map $f:=\Gamma_\Omega:V_A\to V_B$, subject to two positivity conditions inherited from the two primitives:
\begin{equation}\label{eq:two-pos}
f(V_A^+)\subseteq V_B^+\qquad\text{and}\qquad f^*(W_B^+)\subseteq W_A^+ .
\end{equation}
The first asks that $f$ send states to states; the second, now read as a condition on the \emph{same} map $f$ via its transpose, asks that $f$ send effects to effects. Under the no-restriction hypothesis $W_A^+=(V_A^+)^*$, the second condition is automatic: if $f(V_A^+)\subseteq V_B^+$ then for any $\psi\in W_B^+=(V_B^+)^*$ and any $v\in V_A^+$ one has $\langle v,f^*\psi\rangle=\langle f(v),\psi\rangle\ge0$, so $f^*\psi\in(V_A^+)^*=W_A^+$. With restricted effects this argument breaks --- $f^*\psi$ need not land in the smaller cone $W_A^+$ --- and the two conditions in \eqref{eq:two-pos} become genuinely independent constraints.

\subsubsection*{The simplified system}

The two collapses package a system into minimal, $\mathbf{Vec}$-internal data.

\begin{defi}[Simplified GPT system and morphism]
A \emph{system} is a triple $(V_A,V_A^+,W_A^+)$ consisting of a finite-dimensional real vector space $V_A$, a proper state cone $V_A^+\subseteq V_A$, and an effect cone
\[
W_A^+\subseteq(V_A^+)^*\subseteq V_A^*
\]
that \emph{spans} $V_A^*$ and contains a distinguished unit effect $u_A\in W_A^+$. A \emph{morphism} $f:(V_A,V_A^+,W_A^+)\to(V_B,V_B^+,W_B^+)$ is a linear map $f:V_A\to V_B$ that is positive,
\[
f(V_A^+)\subseteq V_B^+,
\]
and whose transpose is $W$-positive,
\[
f^*(W_B^+)\subseteq W_A^+ .
\]
\end{defi}

Every datum now references only $V_A$ and its dual $V_A^*$, both objects of $\mathbf{Vec}$; the abstract effect space and abstract pairing have disappeared. The unit effect $u_A$ and the normalization data $\Omega_A=\{\omega\in V_A^+:\langle\omega,u_A\rangle=1\}$, $\mathcal E_A=[0,u_A]$ are recovered exactly as before, now phrased through evaluation.

\subsubsection*{The dual object}

The simplified data are manifestly symmetric under exchanging the two primitives, and this symmetry is what produces the dual object required for the $1$d theory. Given a system $X=(V_A,V_A^+,W_A^+)$, define its \emph{dual object}
\[
X^\vee:=(V_A^*,\,W_A^+,\,V_A^+)
\]
by swapping the roles of states and effects: the ambient space of $X^\vee$ is $V_A^*$, its state cone is $W_A^+\subseteq V_A^*$, and its effect cone is $V_A^+\subseteq V_A\cong(V_A^*)^*$. The pairing of $X^\vee$ is again evaluation, and the construction is involutive, $(X^\vee)^\vee\cong X$.

For $X^\vee$ to be a legitimate system we must check that its state cone $W_A^+$ generates $V_A^*$ (i.e.\ is generating in the ambient space). This is exactly condition (a) above: $W_A^+$ spans $V_A^*$. Thus the spanning property that survived the failure of no-restriction is precisely what guarantees $X^\vee$ is again a bona fide GPT system. Dually, the effect cone $V_A^+$ of $X^\vee$ spans $V_A$ because $V_A^+$ is generating; and it sits inside $(W_A^+)^*$ because $W_A^+\subseteq(V_A^+)^*$ implies $V_A^+\subseteq(W_A^+)^*$ under the canonical identification $V_A\cong V_A^{**}$.

\subsubsection*{Monoidal adaptation}

Up to this point the reduction has placed a single GPT system entirely inside $\mathbf{Vec}$. The monoidal structure is where the two frameworks genuinely diverge. Recall from the construction of $\square$ that the universal map $\phi_\square:V_A\otimes V_B\to V_{A\square B}$ is injective, yielding the decomposition
\[
V_{A\square B}\;\cong\;(V_A\otimes V_B)\oplus C .
\]
Thus the square product agrees with the ordinary tensor product $\otimes$ of $\mathbf{Vec}$ only on the \emph{locally generated sector} $V_A\otimes V_B$; the complementary global sector $C$ collects degrees of freedom of the composite not accessible to either subsystem alone, and $C=0$ if and only if the theory is locally tomographic.

Even on the local sector the data are not fixed by $\otimes$ alone, because $\mathbf{Vec}$ carries no cones whereas a GPT must specify the composite state cone $V_{A\square B}^+$. Operationally this cone is only bounded, not determined:
\[
\otimes_{\min}\;\subseteq\;V_{A\square B}^+\;\subseteq\;\otimes_{\max},
\]
where $\otimes_{\min}$ is the cone generated by product states $v\square v'$ (the separable cone) and $\otimes_{\max}$ is the cone of all elements positive against every product effect. Summarizing, the square product decomposes into three pieces of data:
\[
\square\;=\;\underbrace{\otimes\big|_{\text{local sector}}}_{\text{the }\mathbf{Vec}\text{ tensor}}\;+\;\underbrace{C}_{\substack{\text{global sector,}\\ \text{zero iff loc.\ tomographic}}}\;+\;\underbrace{[\otimes_{\min},\otimes_{\max}]}_{\text{a choice of composite cone}} .
\]
Only the first piece exists in $\mathbf{Vec}$; the latter two are exactly the additional structure that a GPT carries.

\subsubsection*{GPT versus Vec}

We can now state precisely how the simplified GPT category sits relative to $\mathbf{Vec}$.

\begin{obs}[GPT as a structured enrichment of Vec]\label{obs:gpt-vs-vec}
The simplified GPT category and $\mathbf{Vec}$ are related as follows.
\begin{itemize}
\item \emph{Objects.} An object of $\mathbf{Vec}$ is a bare vector space $V$; a GPT object is $V$ \emph{equipped with} the two cones $(V^+,W^+)$. Forgetting the cones is a functor to $\mathbf{Vec}$ that is the identity on the underlying spaces.
\item \emph{Morphisms.} In $\mathbf{Vec}$ a morphism is an \emph{arbitrary} linear map, so $\operatorname{Hom}_{\mathbf{Vec}}(V_A,V_B)$ is a vector space. In GPT a morphism is a \emph{positive} linear map whose transpose is $W$-positive; positivity is closed under addition and nonnegative scaling but not under subtraction, so $\operatorname{Hom}_{\mathrm{GPT}}(A,B)$ is a \emph{cone} inside the full linear hom-space. Consequently GPT is a \emph{wide} subcategory of $\mathbf{Vec}$ (every object survives the forgetful functor) but a \emph{non-full} one (only the positive maps survive).
\item \emph{Monoidal structure.} $\mathbf{Vec}$ has a single monoidal product $\otimes$ and \emph{is compact closed}: every object is dualizable, with canonical evaluation and coevaluation. The GPT product $\square$ refines $\otimes$ by a global sector $C$ and a choice of composite cone in $[\otimes_{\min},\otimes_{\max}]$, and is \emph{generically not compact closed}: although a dual object $X^\vee$ always exists as a system, the evaluation and coevaluation maps are no longer guaranteed to be \emph{morphisms} (positive maps) for an arbitrary admissible composite cone.
\end{itemize}
\end{obs}

The failure of compact closure flagged in the last point is the genuine departure from $\mathbf{Vec}$ and the crux of what follows: in $\mathbf{Vec}$ dualizability is automatic, whereas in GPT it becomes a contingent positivity property of the chosen cone. The $1$d analysis below makes this precise, locating the entire obstruction in the membership of a single composite element.

\textcolor{red}{[Examples : 1) The classical GPT. 2) The quantum GPT . 3) Fermionic Quantum GPT.  ]}

\section{1-d GPTFTs}

\subsection{Warm-up: 1d topological field theories valued in \texorpdfstring{$\mathbf{Vec}$}{Vec}}

Before turning to GPTs, it is instructive to carry out the one-dimensional construction in full in the simplest possible target, $\mathcal C=\mathbf{Vec}$, the category of finite-dimensional real vector spaces and linear maps. The entire theory will be pinned down by a single integer, and---crucially for the contrast to come---every piece of structure it requires will exist automatically, with no condition left to verify. It is precisely this freedom that fails once $\mathbf{Vec}$ is replaced by the category of GPTs.

\paragraph{The category $\mathbf{Bord}_1$.}
The objects of $\mathbf{Bord}_1$ are closed oriented $0$-manifolds, that is, finite disjoint unions of oriented points. There are two oriented points, $+$ and $-$, and every object is a disjoint union of copies of them. The monoidal product is disjoint union $\sqcup$, the monoidal unit is the empty manifold $\varnothing$, and the symmetry $\sigma$ interchanges two factors of a disjoint union.

The morphisms of $\mathbf{Bord}_1$ are compact oriented $1$-bordisms between such $0$-manifolds, taken up to diffeomorphism relative to the boundary. As a symmetric monoidal category, $\mathbf{Bord}_1$ is generated by three elementary bordisms:
\[
\text{interval}\quad \mathrm{id}:+\to+,
\]
\[
\text{cup}\quad \mathrm{coev}:\varnothing\to(+)\sqcup(-),
\]
\[
\text{cap}\quad \mathrm{ev}:(-)\sqcup(+)\to\varnothing.
\]
The interval is the identity bordism on the point $+$; the cup creates a $+/-$ pair out of nothing, and the cap annihilates a $-/+$ pair. These three generators are not independent: composing a cup with a cap by sliding one boundary point past the other yields a straight interval. Stated as relations between morphisms, this is the pair of \emph{snake} (or \emph{zigzag}) identities
\begin{align}
(\mathrm{id}_+\sqcup\mathrm{ev})\circ(\mathrm{coev}\sqcup\mathrm{id}_+)&=\mathrm{id}_+, \label{eq:snake1}\\
(\mathrm{ev}\sqcup\mathrm{id}_-)\circ(\mathrm{id}_-\sqcup\mathrm{coev})&=\mathrm{id}_-. \label{eq:snake2}
\end{align}
Geometrically, each identity says that a curve which bends down and then back up (an ``S'' shape) can be pulled taut to a straight line, so the two reads of the boundary point agree.

The only other closed $1$-manifold that can be built is the circle. It arises by gluing a cup to a cap after swapping the two intermediate points,
\[
S^1=\mathrm{ev}\circ\sigma\circ\mathrm{coev}\in\mathrm{End}(\varnothing),
\]
and since every closed oriented $1$-manifold is a finite disjoint union of circles, every closed object is $\sqcup_k S^1$ for some $k\ge 0$. Thus the closed invariants of a $1$d TFT are exhausted by the single number assigned to one circle.

\paragraph{The functor and its monoidal rigidity.}
A $1$d TFT valued in $\mathbf{Vec}$ is a symmetric monoidal functor
\[
Z:\mathbf{Bord}_1\to\mathbf{Vec};
\]
equivalently, by the cobordism hypothesis in dimension one, it is the data of a dualizable object of $\mathbf{Vec}$. Symmetric monoidality is extremely restrictive: it forces almost everything from a single choice. Once one fixes the image of the positively oriented point,
\[
Z(+)=V,\qquad \dim_{\mathbb R}V=n,
\]
the remaining values are determined. The empty manifold must go to the monoidal unit, the disjoint union to the tensor product, and the negatively oriented point to the dual:
\[
Z(\varnothing)=\mathbb R,\qquad Z(\sqcup)=\otimes,\qquad Z(-)=V^*.
\]
The interval, being the identity bordism, goes to the identity linear map,
\[
Z(\text{interval})=\mathrm{id}_V.
\]
The only freedom remaining is in the images of the cup and cap, and even these are fixed by the requirement that they furnish the duality between $V$ and $V^*$.

\paragraph{Duality data---canonical and unconditional.}
For any finite-dimensional $V$ there is canonical duality data, defined without any choice. The cap is the evaluation pairing
\[
\mathrm{ev}_V:V^*\otimes V\to\mathbb R,\qquad
\mathrm{ev}_V(\varphi\otimes v)=\varphi(v)=\langle\varphi,v\rangle,
\]
and the cup is the coevaluation
\[
\mathrm{coev}_V:\mathbb R\to V\otimes V^*,\qquad
\mathrm{coev}_V(1)=\sum_a f_a\otimes f^a,
\]
where $\{f_a\}$ is any basis of $V$ and $\{f^a\}$ the dual basis, characterized by $f^a(f_b)=\delta^a_b$. The element $\sum_a f_a\otimes f^a$ is independent of the chosen basis: under the canonical isomorphism $\mathrm{End}(V)\cong V\otimes V^*$ it is exactly $\mathrm{id}_V$, the identity endomorphism written in coordinates. The single computational fact controlling everything below is the resolution of identity,
\begin{align}
v&=\sum_a f^a(v)\,f_a, & v&\in V,\label{eq:resV}\\
\varphi&=\sum_a \varphi(f_a)\,f^a, & \varphi&\in V^*,\label{eq:resVstar}
\end{align}
which simply re-expresses an arbitrary vector or covector in the chosen (co)basis.

\paragraph{The calculation.}
We now verify, on the nose, that the canonical data \eqref{eq:resV}--\eqref{eq:resVstar} satisfy the snakes \eqref{eq:snake1}--\eqref{eq:snake2}. Take a generic $v\in V$ and trace it through the left-hand side of the first snake \eqref{eq:snake1}. First apply $\mathrm{coev}\sqcup\mathrm{id}_+$, which acts as $\mathrm{coev}_V\otimes\mathrm{id}_V$ and tensors $v$ with the coevaluation element:
\[
v\;\longmapsto\;\mathrm{coev}_V(1)\otimes v=\sum_a f_a\otimes f^a\otimes v\ \in\ V\otimes V^*\otimes V.
\]
Now apply $\mathrm{id}_+\sqcup\mathrm{ev}$, which acts as $\mathrm{id}_V\otimes\mathrm{ev}_V$ and contracts the last two tensor factors through the evaluation pairing:
\[
\sum_a f_a\otimes f^a\otimes v\;\longmapsto\;\sum_a f_a\,\mathrm{ev}_V(f^a\otimes v)=\sum_a f_a\,f^a(v)=\sum_a f^a(v)\,f_a.
\]
By the resolution of identity \eqref{eq:resV}, this last sum is exactly $v$. Hence
\[
(\mathrm{id}_+\sqcup\mathrm{ev})\circ(\mathrm{coev}\sqcup\mathrm{id}_+)\,(v)=v\qquad\text{for all }v\in V,
\]
which is the first snake \eqref{eq:snake1}. The second snake \eqref{eq:snake2} follows identically: feeding a generic $\varphi\in V^*$ through $\mathrm{id}_-\sqcup\mathrm{coev}$ and then $\mathrm{ev}\sqcup\mathrm{id}_-$ gives $\sum_a\varphi(f_a)\,f^a$, which equals $\varphi$ by \eqref{eq:resVstar}. No condition is imposed in either case: the snakes are automatic consequences of the resolution of identity, valid for every $V$.

It remains to evaluate the circle. Composing the canonical data as $S^1=\mathrm{ev}\circ\sigma\circ\mathrm{coev}$, we start from $\mathrm{coev}_V(1)=\sum_a f_a\otimes f^a$, swap the two factors with $\sigma$ to get $\sum_a f^a\otimes f_a\in V^*\otimes V$, and contract with $\mathrm{ev}_V$:
\[
Z(S^1)=\mathrm{ev}\circ\sigma\circ\mathrm{coev}\,(1)
=\sum_a \mathrm{ev}_V(f^a\otimes f_a)
=\sum_a f^a(f_a)
=\sum_a 1
=n=\dim_{\mathbb R}V.
\]
This is a special case of the general categorical trace. For any endomorphism $g:V\to V$ one has
\[
\mathrm{ev}\circ\sigma\circ(g\otimes\mathrm{id}_{V^*})\circ\mathrm{coev}
=\sum_a f^a(g\,f_a)
=\operatorname{tr}(g),
\]
so that the circle is precisely $\operatorname{tr}(\mathrm{id}_V)=n$, the categorical dimension of $V$. Combining this with monoidality, a disjoint union of $k$ circles evaluates to
\[
Z(\textstyle\bigsqcup_k S^1)=\big(Z(S^1)\big)^{k}=n^{k},
\]
and, since every closed $1$-manifold is such a union, these $n^k$ exhaust the closed values of the theory.

\paragraph{Classification.}
The computation above shows that the entire theory is determined by the single integer $n$.

\begin{prop}
$1$d topological field theories valued in $\mathbf{Vec}$ are classified, up to isomorphism, by a single nonnegative integer $n=\dim_{\mathbb R}V$, where $V=Z(+)$. Every closed invariant is a power of $n$: a disjoint union of $k$ circles is sent to $n^k$.
\end{prop}

\paragraph{Why it is free.}
It is worth isolating the reason the classification is so rigid, since exactly this freedom will be lost in the GPT setting. The functor $Z$ is pinned by the \emph{single} choice $V=Z(+)$. Once $V$ is fixed, the dual $V^*$, the evaluation $\mathrm{ev}_V$, and the coevaluation $\mathrm{coev}_V$ are \emph{forced}, and they exist \emph{unconditionally}---the pairing $\langle\,\cdot\,,\,\cdot\,\rangle$ and the element $\sum_a f_a\otimes f^a$ are available for any finite-dimensional $V$ with no hypothesis whatsoever. The snake identities then hold \emph{automatically} as a consequence of the resolution of identity, and at no point is any datum tested against a constraint. The whole invariant content of the theory is the integer $\dim_{\mathbb R}V$. This is precisely the freedom that fails in the GPT setting: there the very same cup and cap elements must additionally be \emph{positive}, and that positivity is no longer automatic---it becomes a genuine, contingent condition on the underlying cone, as we develop next.

\subsection{1d GPTFTs: the cone-theoretic departure}

We now specialize the target from $\mathbf{Vec}$ to the category $(\mathrm{GPT},\square)$ of generalized probabilistic systems, in the simplified Vec-like form established above. The categorical data of the bordism category $\mathbf{Bord}_1$ are unchanged; what changes is that every structure map must now be a \emph{positive} morphism of cones, rather than merely a linear map.

\paragraph*{Definition.}
A \emph{$1$d GPTFT} is a symmetric monoidal functor
\[
Z:\mathbf{Bord}_1\to(\mathrm{GPT},\square),
\]
or, equivalently, a dualizable object of $(\mathrm{GPT},\square)$. As in the Vec warm-up, monoidality pins the functor down once a single object is chosen for the positively oriented point. Using the simplified GPT data, the assignment is
\[
Z(+)=X=(V,V^+,W^+),\qquad
Z(-)=X^\vee,\qquad
Z(\varnothing)=\mathbf 1=(\mathbb R,\mathbb R^{\ge0}),
\]
with $Z(\square)=\square$ on disjoint unions and $Z(\text{interval})=\operatorname{id}_X$. Here $V^+\subseteq V$ is the state cone and $W^+\subseteq V^*$ the (spanning) effect cone obtained after discarding the abstract effect space, and the dual object $X^\vee=(V^*,W^+,V^+)$ swaps the two primitives, with $W^+$ spanning precisely so that $X^\vee$ has a generating state cone.

\paragraph*{What transfers for free.}
Everything that was forced in the Vec warm-up survives intact, because the \emph{underlying linear maps are the same}. First, the dual object $X^\vee$ exists by the simplification of the GPT data; no extra input is required to produce it. Second, the cup and cap are carried by exactly the Vec duality maps: the coevaluation is the resolution-of-identity element
\[
\eta(1)=\sum_a f_a\,\square\, f^a\;\in\;V_{X\square X^\vee},
\]
where $\{f_a\}$ is a basis of $V$ and $\{f^a\}$ its dual basis, and the evaluation $\epsilon$ is the canonical pairing $\langle\,\cdot\,,\,\cdot\,\rangle$. Third, the snake identities
\[
(\operatorname{id}_X\square\,\epsilon)\circ(\eta\square\operatorname{id}_X)=\operatorname{id}_X,
\qquad
(\epsilon\square\operatorname{id}_{X^\vee})\circ(\operatorname{id}_{X^\vee}\square\,\eta)=\operatorname{id}_{X^\vee},
\]
are equalities of \emph{linear maps}, established term by term from the resolution of identity
\[
v=\sum_a f^a(v)\,f_a\quad(v\in V),\qquad
\varphi=\sum_a \varphi(f_a)\,f^a\quad(\varphi\in V^*).
\]
They therefore continue to hold automatically, on the nose, \emph{whenever the cup and cap are legal morphisms}; positivity adds no new equation for them to satisfy. Fourth, the value on the circle is forced by the factorization of $\square$ alone, before any positivity is invoked: contracting the two copies of the identity element gives
\[
Z(S^1)=\epsilon\circ\sigma_{X,X^\vee}\circ\eta
=\sum_a f^a(f_a)=\sum_a 1=\dim_{\mathbb R}V .
\]

Operationally, the first snake identity is precisely a teleportation-like transfer protocol. Consider two physicists, Alice and Bob, who share in advance the bipartite resource state $\eta(1)$, with Bob holding the $X$-system and Alice the $X^\vee$-system. Alice receives a positive input state $v\in V_X^+$, so that she now holds two systems, the input $X$ and her share $X^\vee$ of the resource; writing $e:=\epsilon(u_{\mathbf 1})$ for the bipartite effect, she contracts these two systems against $e$. Bob is then left, on his system $X$, with the state obtained by composing $\eta(1)$, the input $v$, and the effect $e$ --- and the snake identity asserts that this state is again $v$. The transfer is exact, and the second snake identity gives the same protocol with the roles of $X$ and $X^\vee$ interchanged. The closed process obtained by removing the external input and output is the trace of this protocol,
\[
Z(S^1)=\operatorname{tr}_{\mathrm{GPT}}(\operatorname{id}_X)=\dim_{\mathbb R}V,
\]
i.e. the categorical dimension of $X$, exactly as in the Vec warm-up where this scalar is the ordinary dimension of the space at the point.

\paragraph*{What changes --- the departure.}
The single thing that is no longer free is the demand that the cup and cap be \emph{morphisms} of $(\mathrm{GPT},\square)$, that is, positive. This is where the whole content of the $1$d theory concentrates. The cup is a state $\eta:\mathbf 1\to X\square X^\vee$, so it must be a \emph{positive bipartite state}; the cap is an effect $\epsilon:X^\vee\square X\to\mathbf 1$, so it must be a \emph{positive bipartite effect}. But by the resolution of identity these are carried by \emph{one and the same} element
\[
g=\sum_a f_a\,\square\, f^a .
\]
Thus $g$ is required to be simultaneously
\begin{itemize}
\item[(a)] a positive state, $g\in V_{X\square X^\vee}^+$ \quad (the \emph{cup}), and
\item[(b)] a positive effect, $g\in\bigl(V_{X^\vee\square X}^+\bigr)^*$ \quad (the \emph{cap}).
\end{itemize}
These two requirements pull in opposite directions. Condition (a) wants the composite state cone to be \emph{large} --- the bigger the cone, the easier it is for $g$ to lie inside it. Condition (b) wants the composite cone to be \emph{small} --- the smaller the cone, the easier it is for $g$ to lie in its dual. Since the composite cone is only pinned to the interval $\otimes_{\min}\subseteq V_{X\square X^\vee}^+\subseteq\otimes_{\max}$, and since min/max duality swaps the two ends of this interval under dualization, the cup and the cap test \emph{one} element against \emph{one} composite cone from inside and from outside at once. In $\mathbf{Vec}$ there is no tension: the same element $g=\operatorname{id}_V$ is unconditionally both a vector and a functional, and dualizability is automatic. In GPT, $g$ must thread the cone from both sides simultaneously, and whether it can is a genuine constraint on the theory.

\paragraph*{Goal statement.}
The problem of classifying $1$d GPTFTs is therefore the following.
\begin{quote}
\emph{Determine for which GPTs --- which state cone $V^+$, and which admissible composite cone in $[\otimes_{\min},\otimes_{\max}]$ --- the identity element $g=\sum_a f_a\square f^a$ is at once a legal cup-state, $g\in V_{X\square X^\vee}^+$, and a legal cap-effect, $g\in(V_{X^\vee\square X}^+)^*$.}
\end{quote}
A $1$d GPTFT exists if and only if some admissible composite cone makes $g$ satisfy both conditions; when it does, the closed value is forced to be $\dim_{\mathbb R}V$. Equivalently: dualizability, which is automatic in $\mathbf{Vec}$, becomes in GPT a \emph{contingent positivity property of the cone}, and that property is concentrated entirely in the membership of the single element $g$.

\paragraph*{Conjectural target.}
\tmp{The following is conjectural and is recorded here as an outlook, not as a theorem.} We conjecture that the invariant which decides the problem is (weak) \emph{self-duality} of the state cone $V^+$. Heuristically, $g$ must lie in the composite cone and in its dual at the same time, and self-duality of $V^+$ is exactly the symmetry that lets a single element occupy both roles. The expected behavior across the standard examples is:
\begin{itemize}
\item \emph{Classical} (orthant) and \emph{quantum} (PSD) cones are self-dual, so they should \emph{pass}. In the classical case $g$ is separable and lies even in the minimal composite $\otimes_{\min}$. In the quantum case $g$ is rescued to the maximally entangled element by a positive but not completely positive transpose twist of the dual leg $X^\vee$, after which it sits inside the genuine composite cone.
\item \emph{Boxworld} (the square cone, with $\ell^\infty\neq\ell^1$) is \emph{not} self-dual, and should give a \emph{no-go}: the cup leg and the cap leg collide quantitatively, an inequality of the schematic form ``$2\le1$''. Because the cone is polyhedral, it admits only a finite group of twists, so there is no continuous rescue analogous to the quantum transpose.
\end{itemize}
This conjectural picture is to be checked against two rigorous results. The first is Barnum--Barrett--Leifer--Wilce \cite{BarnumBarrettLeiferWilce}: weak self-duality together with universal steering implies homogeneity of the state cone. We note that the teleportation/steering structure described above --- the shared resource $\eta(1)$ and the bipartite effect $e=\epsilon(u_{\mathbf 1})$ --- is exactly the cup/cap data appearing in that hypothesis. The second is Koecher--Vinberg \cite{KoecherVinberg}: a homogeneous self-dual cone is precisely the cone of squares of a Euclidean (formally real) Jordan algebra. Together these would place any GPT admitting a $1$d GPTFT ``close to quantum.''

\paragraph*{The subtlety for $C\neq0$.}
A point with no analogue in $\mathbf{Vec}$ arises when the theory is not locally tomographic, i.e. when the global sector $C$ in
\[
V_{X\square X^\vee}\cong (V\otimes V^*)\oplus C
\]
is nonzero. The element $g=\sum_a f_a\square f^a$ lives \emph{entirely} in the local sector $V\otimes V^*$ and has \emph{no} component along $C$. Its positivity, however, is adjudicated inside the \emph{larger} ambient space $(V\otimes V^*)\oplus C$: membership of $g$ in the composite cone and in its dual is decided in all directions, including those of $C$ where $g$ itself is silent and where the composite cone may bend. Keeping $C$ general is therefore not optional --- it is precisely the regime in which a genuinely non-locally-tomographic theory, such as a fermionic GPT, can behave differently from the qudit, even though $g$ never points into $C$.

\paragraph*{Summary.}
A $1$d GPTFT is a dualizable GPT system whose cup/cap positivity data implement the identity on the cone, with value on the circle equal to $\dim_{\mathbb R}V$.

\section{2-d GPTFTs}

\begin{thebibliography}{9}

\bibitem{BarnumBarrettLeiferWilce}
H.~Barnum, J.~Barrett, M.~Leifer, and A.~Wilce,
\emph{Teleportation in general probabilistic theories},
in \emph{Mathematical Foundations of Information Flow},
Proc.\ Sympos.\ Appl.\ Math. (American Mathematical Society, 2012).
% See also H. Barnum, C. P. Gaebler, and A. Wilce, ``Ensemble steering,
% weak self-duality, and the structure of probabilistic theories.''

\bibitem{KoecherVinberg}
M.~Koecher, \emph{The Minnesota Notes on Jordan Algebras and Their Applications},
Lecture Notes in Mathematics \textbf{1710} (Springer, 1999);
E.~B.~Vinberg, \emph{The theory of convex homogeneous cones},
Trans.\ Moscow Math.\ Soc.\ \textbf{12}, 340 (1963).
See also J.~Faraut and A.~Kor\'anyi, \emph{Analysis on Symmetric Cones}
(Oxford University Press, 1994).

\end{thebibliography}

\end{document}
